\subsection{Extraction before the outer challenges}\label{sec:mv:temporal}

\begin{definition}[Pre-challenge registration boundary]\label{def:mv:boundary}
The first outer challenge boundary is the instant
immediately before $\theta$ is sampled.
All commitments to the witness, packed columns,
$m,p,p_{\rm in},c,y'$ that influence the first
discrepancy must be registered before this instant.
All bitness proofs required to bind their
binary-coordinate interpretations are completed
before this boundary.
The second boundary is immediately before $z$
is sampled. The post-$\theta$ variables
$t,\rho,\sigma$ and their required bitness proofs
are completed before this second boundary.
The \emph{registration prologue} at either boundary
consists of the appropriate bit commitments
followed by sequential one-bit OR proof repetitions,
with no outer random challenge interleaved
inside that prologue.
\end{definition}

\begin{lemma}[A local reset occurs before the protected challenge]\label{lem:mv:temporal-local}
If the entire OR prologue for a statement
$C,S$ is completed before an outer challenge
is sampled, then the two-run forking experiment
of Theorem~\ref{thm:mv:fork-average} can be
performed using only states and verifier
challenge bits from \emph{before} that
outer challenge.
\end{lemma}
\begin{proof}
The prologue begins with fixed public
statement and a prover state prior to
the protected outer challenge.
Both forked runs replay the prologue
from this saved state and use only
new independent \emph{inner} verifier bits.
By the defined schedule, neither run
needs or receives the future outer
challenge. The extractor's messages,
including the second branch challenge
that creates the fork, are therefore
all generated before the outer challenge.
\end{proof}

\begin{theorem}[Stand-alone early extraction with quantitative success]\label{thm:mv:standalone-early}
Fix one registration OR statement and a
resettable polynomial-time prover in a
$\kappa$-copy prologue that ends before
an outer challenge. If its probability
of passing this prologue against
independent uniform challenge bits is
$\varepsilon$, then a two-run
extractor \emph{before that outer challenge}
produces a valid permitted opening of the
statement with probability at least
\[
 \boxed{\max\{0,\varepsilon(\varepsilon-2^{-\kappa})\}.}
\]
Its two executions use no outer challenge
value and therefore do not condition on
one.
\end{theorem}
\begin{proof}
Lemma~\ref{lem:mv:temporal-local}
places the fork entirely in the
pre-challenge prologue.
Theorem~\ref{thm:mv:fork-average}
gives the probability of two
distinct accepting challenge strings,
and Theorem~\ref{thm:mv:two-run-extract}
computes an allowed witness from the
resulting fork. All actions happen
before the protected outer challenge
is sampled.
\end{proof}

\begin{lemma}[At most one successful challenge path without a fork]\label{lem:mv:single-leaf}
For a fixed deterministic prover state and
random tape, if no pair of distinct
accepting challenge strings yields
a valid fork, then at most one of
the $2^\kappa$ challenge strings can
be accepting; its acceptance probability
is therefore at most $2^{-\kappa}$.
\end{lemma}
\begin{proof}
If two distinct accepting challenge
strings existed, Lemma~\ref{lem:mv:earliest}
would yield a pair of accepting
responses to one first message
under different challenge bits.
That is a valid fork, contrary to
hypothesis. Thus the accepting set
has cardinality at most one.
Its uniform probability is
at most $1/2^\kappa$.
\end{proof}

\begin{lemma}[False parity plus a fork solves setup discrete log]\label{lem:mv:false-parity-fork}
Suppose the registered bits of an
aggregate commitment are fixed and
the extractor knows valid binary
openings $(b_i,r_i)$.
Let $S=\sum_ib_i$ be its integer
sum, and assume the claimed parity
target is false for these bits.
If an accepting $\kappa$-copy
aggregate parity OR proof yields
a fork, one can compute
$\alpha=\log_g h$ in polynomial time.
\end{lemma}
\begin{proof}
Theorem~\ref{thm:mv:two-run-extract}
extracts an opening $(t,R)$ with
$t$ in the parity-compatible candidate
set. The aggregate commitment also
has the known true opening
$(S,\sum_ir_i)$ by
Lemma~\ref{lem:realize:aggregate}.
Because the target parity is false
and $P$ exceeds the row length,
Lemma~\ref{lem:realize:parity-equivocation}
shows that $t$ and $S$ are distinct
modulo $P$. The explicit formula of
Theorem~\ref{thm:realize:parity-binding}
then recovers the setup logarithm.
\end{proof}

\begin{theorem}[Quantitative soundness for a false parity with known true opening]\label{thm:mv:false-parity-bound}
Suppose the correct binary openings
of a fixed aggregate commitment are
known to the reduction, but their
parity disagrees with the public
permitted candidate set. Let
$\varepsilon$ be the success
probability of a causal prover
against $\kappa$ sequential honest
one-bit challenges.
There is a two-run discrete-log
reduction whose success probability
is at least
\[
 \boxed{
 \max\{0,\varepsilon(\varepsilon-2^{-\kappa})\}.}
\]
Consequently, if that reduction's
discrete-log success probability is
at most $\delta_{\rm DL}$, then
\[
 \boxed{\varepsilon\le
 2^{-\kappa}+\sqrt{\delta_{\rm DL}}.}
\]
\end{theorem}
\begin{proof}
Apply the two-run forking bound
Theorem~\ref{thm:mv:fork-average}.
For every successful fork,
Lemma~\ref{lem:mv:false-parity-fork}
computes the forbidden logarithm,
giving the first inequality.
Write $a=2^{-\kappa}$.
If $\varepsilon\le a$, the
second bound holds trivially.
Otherwise
$\varepsilon(\varepsilon-a)
\le\delta_{\rm DL}$.
Since $\varepsilon\ge\varepsilon-a$,
we have
$(\varepsilon-a)^2\le
\delta_{\rm DL}$.
Taking nonnegative square roots
yields $\varepsilon\le a+
\sqrt{\delta_{\rm DL}}$.
\end{proof}

\begin{lemma}[Different first messages do not imply extractability]\label{lem:mv:fork-needs-same-first}
Two accepting OR mini-proof transcripts
with the same statement but different
first-message vectors are not sufficient
to apply the special-soundness formula
of Theorem~\ref{thm:realize:or-extract}.
\end{lemma}
\begin{proof}
Special soundness explicitly assumes
the same $T$ in both transcripts.
Indeed, for any public OR statement
and either challenge bit, the
witness-free algorithm of
Definition~\ref{def:mv:fixed-sim}
generates accepting transcripts.
Two independent runs generally
produce different first messages
yet require no witness to generate.
Thus the mere existence of two
accepting transcripts with
different first messages supplies
none of the algebraic cancellation
used by the extraction proof.
\end{proof}

\begin{definition}[Global temporal coupling obligation]\label{def:mv:global-coupling}
The \emph{global extraction obligation}
is to build an efficient extractor for
the actual adaptive interaction that:
\begin{enumerate}[label=\textup{(\roman*)}]
\item assigns a unique extracted binary
value to \emph{all} pre-$\theta$
committed field variables before
sampling the real $\theta$;
\item similarly extracts all required
post-$\theta$, pre-$z$ values before
the real $z$;
\item preserves a \emph{single}
consistent accepted visible transcript,
including its already sampled outer
challenge and authenticated openings;
\item keeps the real independent
outer challenge laws intact in the
coupled ideal run; and
\item incurs an explicitly bounded
bad-event probability efficient enough
to combine with the algebraic
soundness theorem of Chapter~9.
\end{enumerate}
This is the stronger contract of
Definition~\ref{def:realize:early-extract}
with a fully instantiated, efficient
multi-statement extractor.
\end{definition}

\begin{proposition}[Local extraction does not yet prove global coupling]\label{prop:mv:coupling-open}
The extraction theorem
Theorem~\ref{thm:mv:standalone-early}
does not by itself discharge the
global obligation of
Definition~\ref{def:mv:global-coupling}.
\end{proposition}
\begin{proof}
The stand-alone theorem concerns
one fixed statement, one saved
initial prover state and two
replayed \emph{inner} challenge
sequences. It yields a witness
on the fork event but says
nothing about selecting one
accepted continuation that
simultaneously matches all
other registered bit statements,
their dependent affine rows,
and the later public history.
The global obligation additionally
requires consistent extraction
of all statements, a coupling of
the final real accepted transcript,
and preservation of the independent
outer challenges. None of these
conclusions is among the hypotheses
or statements proved for one
isolated OR prologue. A separate
multi-stage extractor is therefore
required; no impossibility is claimed.
\end{proof}
